\documentclass[10pt,a4paper]{article}

\usepackage[utf8]{inputenc}
\usepackage[T2A]{fontenc}
\usepackage[russian,english]{babel}
\usepackage[a4paper,margin=1.08cm]{geometry}
\usepackage{enumitem}
\usepackage{titlesec}
\usepackage{xcolor}
\usepackage[colorlinks=true,urlcolor=blue,linkcolor=black]{hyperref}
\usepackage{cmap}

\urlstyle{same}
\setlist[itemize]{leftmargin=1.02em,itemsep=0.06em,topsep=0.06em,parsep=0em,partopsep=0em}
\titleformat{\section}{\normalsize\bfseries}{}{0em}{}
\titlespacing*{\section}{0pt}{0.24em}{0.11em}
\pagestyle{empty}
\setlength{\parindent}{0pt}
\setlength{\parskip}{0pt}
\emergencystretch=2em

\begin{document}
\small

\begin{center}
{\LARGE \textbf{Artem Sidnev}}\\[0.20em]
{\normalsize Инженер-программист (бэкенд)}\\[0.35em]
{\footnotesize
\href{mailto:a.sidnevart@gmail.com}{a.sidnevart@gmail.com}
\enspace\textbar\enspace
\href{https://t.me/sidnevart}{Telegram}
\enspace\textbar\enspace
\href{https://www.linkedin.com/in/artem-sidnev}{LinkedIn}
\enspace\textbar\enspace
\href{https://github.com/sidnevart}{GitHub}
\enspace\textbar\enspace
\href{https://habr.com/ru/users/sidnevart/}{Habr}
\enspace\textbar\enspace
\href{https://x.com/ArtemkaWeb3}{X/Twitter}
}
\end{center}

\vspace{0.06em}

\section*{Кратко}
Инженер-программист с почти 3 годами опыта в бэкенде, платформенных решениях, автоматизации и AI-инструментах. С января 2025 года работаю в T-Bank над платформой таргетирования кэшбэков для партнёрских предложений. С августа 2024 по август 2025 работал над платформой аналитики коммерческой недвижимости для крупного агентства недвижимости; фриланс-проекты в бэкенде и автоматизации веду с июня 2023 года.

\section*{Опыт}

\textbf{Инженер-программист (бэкенд) \textbar{} T-Bank} \hfill \textit{янв 2025 -- н.в.}
\begin{itemize}
  \item Разрабатываю бэкенд для платформы таргетирования кэшбэков, которая обрабатывает аудитории от партнёров и помогает отправлять релевантные офферы нужным клиентам.
  \item Реализовал таргетирование по тратам и доходам клиентов на потоке около 30 млн транзакций в день и 1.5 лет исторических данных; годовой бизнес-эффект направления составляет порядка \$600 тыс.+.
  \item Участвовал в переносе процессов сборки аудиторий во внутреннюю аналитическую среду ClickHouse as a Service, что ускорило критичный сценарий примерно в 5 раз и позволило убрать нестандартные процессы из устаревшего контура.
  \item Автоматизировал подстановку регионов партнёров при настройке аудиторий, убрав 10--15 минут с одного запуска и до 40--60 часов ручной работы в месяц.
  \item Добавил раннюю проверку в критичных сценариях таргетинга, сократил диагностику типовой проблемы на 20--40 минут и ускорил основной CI-пайплайн на 5--7 минут, то есть примерно на 15--20\%.
  \item Решил серьёзный инцидент в сервисе документооборота, устранив репутационные риски для банка и предотвратив потенциальную потерю денежных средств.
  \item Внедрял AI во внутренние прикладные процессы: AI-проверку кода в GitLab, AI-задачу в CI для автоматического рефакторинга и подготовки черновиков merge request, n8n-автоматизацию для менеджеров и RAG-агента для адаптации стажёров.
\end{itemize}

\textbf{Платформа аналитики коммерческой недвижимости \textbar{} крупное агентство недвижимости} \hfill \textit{авг 2024 -- авг 2025}
\begin{itemize}
  \item Работал над платформой аналитики коммерческой недвижимости, которая объединяла данные аукционов, CIAN, геокодинг, Telegram-бота, административную панель и каталог объектов.
  \item Платформа позволяла риэлторам агентства за 10 минут анализировать экономику объекта: прибыль по продаже и аренде, ключевые метрики, денежный поток, доходность и срок окупаемости вместо ручного мониторинга рынка.
  \item Реализовал быстрый анализ рынка по конкретному району за 7--10 минут и удобную навигацию по объектам через административную панель и структурированный каталог.
\end{itemize}

\textbf{Фриланс-проекты в бэкенде и автоматизации} \hfill \textit{июн 2023 -- н.в.}
\begin{itemize}
  \item Работал над фриланс-проектами в бэкенде и автоматизации, накапливая практический опыт в сервисах, интеграциях и автоматизации процессов до и параллельно с продуктовой работой.
\end{itemize}

\section*{Открытый проект}

\textbf{ARC \textbar{} платформа с локальным хранением данных для работы с AI-агентами}
\begin{itemize}
  \item Создал среду с локальным хранением данных для безопасной и воспроизводимой работы с AI-агентами в кодовых проектах.
  \item Объединил в одном продукте командную среду на Go, настольное приложение на Wails, адаптеры к нескольким AI-провайдерам, локальную память проекта и систему артефактов запусков.
  \item Настроил управление контекстом, памятью и правилами выполнения, чтобы снизить стоимость AI-запросов, сделать результат более предсказуемым и упростить проверку каждого шага; сократил контекстную нагрузку на 48\% без ухудшения качества ответа.
  \item Добавил встроенные диаграммы, мини-приложения, демонстрации и симуляции прямо в ответы агента, чтобы результат было проще понять, проверить и использовать в работе.
\end{itemize}

\section*{Стек}
Java, Kotlin, Go, Python, TypeScript, Spring Boot, PostgreSQL, ClickHouse, Redis, Kafka, GitLab CI, n8n, RAG, AI-агенты, внутренние платформы, автоматизация процессов, аналитические сервисы, CRM, партнёрские интеграции, бэкенд-разработка и системы с интенсивной работой с данными.

\section*{Образование}
\textbf{Canisius University} \textbar{} Bachelor of Computer Science\\
\textbf{Дополнительно:} Harvard CS50, MIT OpenCourseWare Algorithms, freeCodeCamp Java, freeCodeCamp Spring

\end{document}